\chapter{Unit 5: Disaster Management and Hazard Mitigation}

\section{Section A: Multiple Choice Practice Questions (60 Curated Problems)}

\noindent\textbf{Q1. Which of the following is a natural disaster?} \hfill [Tag: Natural disasters | Easy]
\begin{enumerate}[label=(\alph*),topsep=2pt,itemsep=1pt,partopsep=0pt,parsep=0pt]
    \item {Earthquake}
    \item {Oil spill}
    \item {Factory fire}
    \item {Train collision}
\end{enumerate}
\vspace{2mm}

\noindent\textbf{Q2. What is a flood?} \hfill [Tag: Water-related disasters: floods and drought | Easy]
\begin{enumerate}[label=(\alph*),topsep=2pt,itemsep=1pt,partopsep=0pt,parsep=0pt]
    \item {Long period without rainfall}
    \item {Overflow of water onto land}
    \item {Sudden movement of rocks}
    \item {Strong rotation of winds}
\end{enumerate}
\vspace{2mm}

\noindent\textbf{Q3. Which condition is most directly associated with drought?} \hfill [Tag: Water-related disasters: floods and drought | Easy]
\begin{enumerate}[label=(\alph*),topsep=2pt,itemsep=1pt,partopsep=0pt,parsep=0pt]
    \item {Prolonged lack of rainfall}
    \item {Sudden release of lava}
    \item {Continuous rise in tides}
    \item {Rapid movement of snow}
\end{enumerate}
\vspace{2mm}

\noindent\textbf{Q4. Where do tropical cyclones generally develop?} \hfill [Tag: Air-related disasters: cyclones and storms | Easy]
\begin{enumerate}[label=(\alph*),topsep=2pt,itemsep=1pt,partopsep=0pt,parsep=0pt]
    \item {Over warm ocean waters}
    \item {On permanently frozen land}
    \item {Below the Earth's surface}
    \item {Inside deep river valleys}
\end{enumerate}
\vspace{2mm}

\noindent\textbf{Q5. Which action is safest during a severe storm?} \hfill [Tag: Air-related disasters: cyclones and storms | Easy]
\begin{enumerate}[label=(\alph*),topsep=2pt,itemsep=1pt,partopsep=0pt,parsep=0pt]
    \item {Walk through flooded streets}
    \item {Stand beneath a tall tree}
    \item {Remain in an open field}
    \item {Stay inside a secure building}
\end{enumerate}
\vspace{2mm}

\noindent\textbf{Q6. Which instrument records earthquake vibrations?} \hfill [Tag: Earth-related disasters: earthquakes, landslides, avalanches and volcanic eruptions | Easy]
\begin{enumerate}[label=(\alph*),topsep=2pt,itemsep=1pt,partopsep=0pt,parsep=0pt]
    \item {Rain gauge}
    \item {Anemometer}
    \item {Barometer}
    \item {Seismograph}
\end{enumerate}
\vspace{2mm}

\noindent\textbf{Q7. What is a landslide?} \hfill [Tag: Earth-related disasters: earthquakes, landslides, avalanches and volcanic eruptions | Easy]
\begin{enumerate}[label=(\alph*),topsep=2pt,itemsep=1pt,partopsep=0pt,parsep=0pt]
    \item {Upward movement of ocean water}
    \item {Horizontal movement of sea ice}
    \item {Circular movement of strong wind}
    \item {Downward movement of rock and soil}
\end{enumerate}
\vspace{2mm}

\noindent\textbf{Q8. An avalanche mainly involves the rapid movement of which material?} \hfill [Tag: Earth-related disasters: earthquakes, landslides, avalanches and volcanic eruptions | Easy]
\begin{enumerate}[label=(\alph*),topsep=2pt,itemsep=1pt,partopsep=0pt,parsep=0pt]
    \item {Lava and ash}
    \item {Water and mud}
    \item {Sand and dust}
    \item {Snow and ice}
\end{enumerate}
\vspace{2mm}

\noindent\textbf{Q9. What may be released during a volcanic eruption?} \hfill [Tag: Earth-related disasters: earthquakes, landslides, avalanches and volcanic eruptions | Easy]
\begin{enumerate}[label=(\alph*),topsep=2pt,itemsep=1pt,partopsep=0pt,parsep=0pt]
    \item {Snow, ice, and hail}
    \item {Lava, ash, and gases}
    \item {Sand, soil, and roots}
    \item {Rain, mist, and fog}
\end{enumerate}
\vspace{2mm}

\noindent\textbf{Q10. Which event is classified as a manmade disaster?} \hfill [Tag: Manmade disasters | Easy]
\begin{enumerate}[label=(\alph*),topsep=2pt,itemsep=1pt,partopsep=0pt,parsep=0pt]
    \item {Volcanic eruption}
    \item {Tropical cyclone}
    \item {Industrial gas leak}
    \item {Natural avalanche}
\end{enumerate}
\vspace{2mm}

\noindent\textbf{Q11. What is the main hazardous release associated with a nuclear disaster?} \hfill [Tag: Nuclear disasters | Easy]
\begin{enumerate}[label=(\alph*),topsep=2pt,itemsep=1pt,partopsep=0pt,parsep=0pt]
    \item {Organic compost}
    \item {Volcanic magma}
    \item {Rainwater runoff}
    \item {Radioactive material}
\end{enumerate}
\vspace{2mm}

\noindent\textbf{Q12. Which incident is an example of a chemical disaster?} \hfill [Tag: Chemical disasters | Easy]
\begin{enumerate}[label=(\alph*),topsep=2pt,itemsep=1pt,partopsep=0pt,parsep=0pt]
    \item {Collapse of a snow slope}
    \item {Leak of toxic gas}
    \item {Shaking of the ground}
    \item {Overflow of a river}
\end{enumerate}
\vspace{2mm}

\noindent\textbf{Q13. Which event is considered a biological disaster?} \hfill [Tag: Biological disasters | Easy]
\begin{enumerate}[label=(\alph*),topsep=2pt,itemsep=1pt,partopsep=0pt,parsep=0pt]
    \item {Widespread infectious disease outbreak}
    \item {Powerful coastal cyclone}
    \item {Major railway collision}
    \item {Large offshore oil spill}
\end{enumerate}
\vspace{2mm}

\noindent\textbf{Q14. Which event is a transport accident?} \hfill [Tag: Transport accidents | Easy]
\begin{enumerate}[label=(\alph*),topsep=2pt,itemsep=1pt,partopsep=0pt,parsep=0pt]
    \item {Volcano releasing ash}
    \item {Drought damaging crops}
    \item {River flooding its banks}
    \item {Passenger train derailment}
\end{enumerate}
\vspace{2mm}

\noindent\textbf{Q15. What is the primary purpose of a national disaster management framework?} \hfill [Tag: National disaster management framework | Easy]
\begin{enumerate}[label=(\alph*),topsep=2pt,itemsep=1pt,partopsep=0pt,parsep=0pt]
    \item {Manage international sports events}
    \item {Regulate daily weather changes}
    \item {Coordinate disaster risk and response}
    \item {Set prices for farm products}
\end{enumerate}
\vspace{2mm}

\noindent\textbf{Q16. In India, which body lays down national policies and guidelines for disaster management?} \hfill [Tag: National disaster management framework | Easy]
\begin{enumerate}[label=(\alph*),topsep=2pt,itemsep=1pt,partopsep=0pt,parsep=0pt]
    \item {National Highways Authority}
    \item {University Grants Commission}
    \item {National Disaster Management Authority}
    \item {Election Commission of India}
\end{enumerate}
\vspace{2mm}

\noindent\textbf{Q17. Which task is commonly performed by governmental agencies during a disaster?} \hfill [Tag: Role of governmental agencies in disaster response | Easy]
\begin{enumerate}[label=(\alph*),topsep=2pt,itemsep=1pt,partopsep=0pt,parsep=0pt]
    \item {Promoting holiday travel}
    \item {Conducting entertainment programs}
    \item {Advertising commercial products}
    \item {Organizing evacuation and rescue}
\end{enumerate}
\vspace{2mm}

\noindent\textbf{Q18. How do non-governmental organizations commonly help disaster-affected people?} \hfill [Tag: Role of NGOs in disaster response | Easy]
\begin{enumerate}[label=(\alph*),topsep=2pt,itemsep=1pt,partopsep=0pt,parsep=0pt]
    \item {Operating all military installations}
    \item {Distributing food and essential supplies}
    \item {Changing official weather records}
    \item {Issuing national currency and coins}
\end{enumerate}
\vspace{2mm}

\noindent\textbf{Q19. Why are community-based organizations valuable during disaster response?} \hfill [Tag: Role of community-based organizations in disaster response | Easy]
\begin{enumerate}[label=(\alph*),topsep=2pt,itemsep=1pt,partopsep=0pt,parsep=0pt]
    \item {They predict the exact time of disasters}
    \item {They understand local needs and resources}
    \item {They replace all professional rescue teams}
    \item {They control every national emergency agency}
\end{enumerate}
\vspace{2mm}

\noindent\textbf{Q20. What is an important role of the media during a disaster?} \hfill [Tag: Role of media in disaster response | Easy]
\begin{enumerate}[label=(\alph*),topsep=2pt,itemsep=1pt,partopsep=0pt,parsep=0pt]
    \item {Blocking communication between agencies}
    \item {Hiding official evacuation information}
    \item {Sharing verified warnings and instructions}
    \item {Spreading unconfirmed alarming rumors}
\end{enumerate}
\vspace{2mm}

\noindent\textbf{Q21. A district is exposed to earthquakes, but most buildings follow seismic codes and residents regularly conduct evacuation drills. Which statement best describes the district's disaster risk?} \hfill [Tag: Natural disasters | Medium]
\begin{enumerate}[label=(\alph*),topsep=2pt,itemsep=1pt,partopsep=0pt,parsep=0pt]
    \item {Risk is unchanged because only the hazard determines disaster losses}
    \item {Risk is increased because preparedness activities create public anxiety}
    \item {Risk is reduced because vulnerability and lack of preparedness are lower}
    \item {Risk is eliminated because earthquakes can no longer affect the district}
\end{enumerate}
\vspace{2mm}

\noindent\textbf{Q22. A town receives an intense rainfall warning after several days of rain, and the soil is already saturated. Which immediate action would most effectively reduce loss of life from a possible flash flood?} \hfill [Tag: Water-related disasters: floods and drought | Medium]
\begin{enumerate}[label=(\alph*),topsep=2pt,itemsep=1pt,partopsep=0pt,parsep=0pt]
    \item {Move residents from low-lying areas to designated higher shelters}
    \item {Distribute drought-resistant seeds to farmers near the floodplain}
    \item {Begin constructing a new concrete dam across the main river}
    \item {Deepen every drainage channel before the rainfall begins that evening}
\end{enumerate}
\vspace{2mm}

\noindent\textbf{Q23. A drought-prone farming region relies heavily on a rapidly declining groundwater supply. Which measure offers the most sustainable reduction in drought vulnerability?} \hfill [Tag: Water-related disasters: floods and drought | Medium]
\begin{enumerate}[label=(\alph*),topsep=2pt,itemsep=1pt,partopsep=0pt,parsep=0pt]
    \item {Replace food crops entirely with water-intensive cash crops}
    \item {Transport water by tanker as the permanent regional water policy}
    \item {Pump deeper groundwater during every dry season}
    \item {Adopt rainwater harvesting and water-efficient irrigation}
\end{enumerate}
\vspace{2mm}

\noindent\textbf{Q24. A coastal village receives a cyclone warning that predicts a severe storm surge. Which location is safest for evacuation?} \hfill [Tag: Air-related disasters: cyclones and storms | Medium]
\begin{enumerate}[label=(\alph*),topsep=2pt,itemsep=1pt,partopsep=0pt,parsep=0pt]
    \item {A basement below a strong building near the shoreline}
    \item {A fishing boat secured inside the village harbor}
    \item {A designated cyclone shelter on elevated inland ground}
    \item {A concrete house located beside the beach}
\end{enumerate}
\vspace{2mm}

\noindent\textbf{Q25. During a severe thunderstorm, a family is inside a well-built house. Which action is most appropriate?} \hfill [Tag: Air-related disasters: cyclones and storms | Medium]
\begin{enumerate}[label=(\alph*),topsep=2pt,itemsep=1pt,partopsep=0pt,parsep=0pt]
    \item {Stand on the balcony to observe changes in wind direction}
    \item {Shelter beneath the tallest nearby tree until lightning stops}
    \item {Stay indoors and avoid windows and electrical appliances}
    \item {Use wired electrical devices to monitor all available weather reports continuously}
\end{enumerate}
\vspace{2mm}

\noindent\textbf{Q26. A student is inside a classroom when strong earthquake shaking begins. What should the student do first?} \hfill [Tag: Earth-related disasters: earthquakes, landslides, avalanches and volcanic eruptions | Medium]
\begin{enumerate}[label=(\alph*),topsep=2pt,itemsep=1pt,partopsep=0pt,parsep=0pt]
    \item {Run immediately toward the staircase during the shaking}
    \item {Stand beside a window to watch for falling objects}
    \item {Drop, cover, and hold on beneath a sturdy desk}
    \item {Use the elevator to reach the ground floor as quickly as possible}
\end{enumerate}
\vspace{2mm}

\noindent\textbf{Q27. After several days of heavy rainfall, cracks appear across a steep hillside above a settlement. What is the most appropriate response?} \hfill [Tag: Earth-related disasters: earthquakes, landslides, avalanches and volcanic eruptions | Medium]
\begin{enumerate}[label=(\alph*),topsep=2pt,itemsep=1pt,partopsep=0pt,parsep=0pt]
    \item {Fill the visible cracks with loose soil and continue normal activity}
    \item {Evacuate the exposed area and have the slope assessed}
    \item {Remove vegetation from the slope to reduce its total weight}
    \item {Direct additional runoff toward the cracked area to identify weak sections}
\end{enumerate}
\vspace{2mm}

\noindent\textbf{Q28. Authorities detect increasing volcanic tremors, gas emissions, and ground swelling near an active volcano. What do these observations most reasonably indicate?} \hfill [Tag: Earth-related disasters: earthquakes, landslides, avalanches and volcanic eruptions | Medium]
\begin{enumerate}[label=(\alph*),topsep=2pt,itemsep=1pt,partopsep=0pt,parsep=0pt]
    \item {A drought is reducing pressure beneath the volcano}
    \item {Magma may be moving upward and eruption risk is increasing}
    \item {The volcano has become permanently inactive}
    \item {The warning signs confirm the exact time and size of the next eruption}
\end{enumerate}
\vspace{2mm}

\noindent\textbf{Q29. A large fire starts in an overcrowded public hall where an emergency exit has been locked. Which preparedness failure contributes most directly to the likely casualties?} \hfill [Tag: Manmade disasters | Medium]
\begin{enumerate}[label=(\alph*),topsep=2pt,itemsep=1pt,partopsep=0pt,parsep=0pt]
    \item {Failure to schedule events during daylight hours}
    \item {Failure to maintain accessible evacuation routes}
    \item {Failure to install decorative signs at the entrance}
    \item {Failure to provide detailed printed information about the building's construction history}
\end{enumerate}
\vspace{2mm}

\noindent\textbf{Q30. Following an accidental release of radioactive material, authorities instruct residents in a nearby area to shelter indoors. Which action best follows this instruction?} \hfill [Tag: Nuclear disasters | Medium]
\begin{enumerate}[label=(\alph*),topsep=2pt,itemsep=1pt,partopsep=0pt,parsep=0pt]
    \item {Stand outside to determine the direction of the wind}
    \item {Drink iodine solution prepared from household disinfectant}
    \item {Travel immediately through the affected zone without waiting for an approved evacuation route}
    \item {Close doors and windows and switch off outside-air ventilation}
\end{enumerate}
\vspace{2mm}

\noindent\textbf{Q31. Why should potassium iodide tablets be taken after a nuclear emergency only when public health authorities recommend them?} \hfill [Tag: Nuclear disasters | Medium]
\begin{enumerate}[label=(\alph*),topsep=2pt,itemsep=1pt,partopsep=0pt,parsep=0pt]
    \item {They protect mainly the thyroid from radioactive iodine and require proper timing}
    \item {They permanently neutralize contaminated soil, food, water, clothing, and equipment}
    \item {They remove every radioactive substance from the entire body}
    \item {They prevent external radiation from penetrating buildings}
\end{enumerate}
\vspace{2mm}

\noindent\textbf{Q32. A tanker releases a toxic gas that is moving downwind toward a settlement. If evacuation is ordered, which route should residents generally use?} \hfill [Tag: Chemical disasters | Medium]
\begin{enumerate}[label=(\alph*),topsep=2pt,itemsep=1pt,partopsep=0pt,parsep=0pt]
    \item {Move crosswind toward an identified safe area}
    \item {Move toward the tanker to obtain accurate information}
    \item {Move directly downwind to remain ahead of the gas}
    \item {Follow the lowest drainage channel because toxic gases always rise rapidly}
\end{enumerate}
\vspace{2mm}

\noindent\textbf{Q33. An unknown chemical has spilled onto a worker's clothing in a factory. What should trained responders generally prioritize after moving the worker away from the source?} \hfill [Tag: Chemical disasters | Medium]
\begin{enumerate}[label=(\alph*),topsep=2pt,itemsep=1pt,partopsep=0pt,parsep=0pt]
    \item {Allow the chemical to dry before examining the affected skin}
    \item {Apply a neutralizing chemical without identifying the spilled substance}
    \item {Send the worker home while preserving the contaminated clothing for later laboratory analysis}
    \item {Remove contaminated clothing and begin appropriate decontamination}
\end{enumerate}
\vspace{2mm}

\noindent\textbf{Q34. Several patients from the same neighborhood develop an unusual infectious illness. Which early public health action is most useful for controlling a possible biological disaster?} \hfill [Tag: Biological disasters | Medium]
\begin{enumerate}[label=(\alph*),topsep=2pt,itemsep=1pt,partopsep=0pt,parsep=0pt]
    \item {Close all hospitals to prevent additional patients from arriving}
    \item {Distribute the same antibiotic to the whole population before identifying the disease}
    \item {Wait for every laboratory result before reporting the cluster}
    \item {Identify cases, trace contacts, and strengthen disease surveillance}
\end{enumerate}
\vspace{2mm}

\noindent\textbf{Q35. A truck carrying hazardous chemicals overturns on a highway. Bystanders notice a leaking container but no visible flames. What should they do?} \hfill [Tag: Transport accidents | Medium]
\begin{enumerate}[label=(\alph*),topsep=2pt,itemsep=1pt,partopsep=0pt,parsep=0pt]
    \item {Move the damaged containers by hand so that normal traffic can resume}
    \item {Wash the leaking material into the roadside drainage system}
    \item {Keep away, move upwind, and notify emergency services}
    \item {Approach the truck and identify the substance by smell}
\end{enumerate}
\vspace{2mm}

\noindent\textbf{Q36. Under India's national disaster management framework, which approach best reflects the shift from relief-centered management to comprehensive disaster management?} \hfill [Tag: National disaster management framework | Medium]
\begin{enumerate}[label=(\alph*),topsep=2pt,itemsep=1pt,partopsep=0pt,parsep=0pt]
    \item {Concentrating resources only on compensation after each disaster}
    \item {Assigning all disaster functions exclusively to the central government}
    \item {Integrating prevention, mitigation, preparedness, response, and recovery}
    \item {Replacing local emergency plans with one identical national procedure for every hazard}
\end{enumerate}
\vspace{2mm}

\noindent\textbf{Q37. After a major flood affects several districts, which task is most appropriately led by governmental agencies?} \hfill [Tag: Role of governmental agencies in disaster response | Medium]
\begin{enumerate}[label=(\alph*),topsep=2pt,itemsep=1pt,partopsep=0pt,parsep=0pt]
    \item {Transfer all responsibility for rescue and public order to spontaneous volunteer groups}
    \item {Allow each responding department to operate without a shared command structure}
    \item {Publish unverified casualty estimates before field reports are received}
    \item {Coordinate evacuation, public services, relief distribution, and interagency support}
\end{enumerate}
\vspace{2mm}

\noindent\textbf{Q38. An NGO arrives in a cyclone-affected area where government teams are already distributing food. What should the NGO do first to maximize its usefulness?} \hfill [Tag: Role of NGOs in disaster response | Medium]
\begin{enumerate}[label=(\alph*),topsep=2pt,itemsep=1pt,partopsep=0pt,parsep=0pt]
    \item {Coordinate with authorities and assess unmet community needs}
    \item {Begin an independent distribution at the same locations}
    \item {Select beneficiaries according to the NGO's publicity requirements}
    \item {Import unfamiliar supplies in large quantities without checking local needs, customs, storage, or transport capacity}
\end{enumerate}
\vspace{2mm}

\noindent\textbf{Q39. Why can a trained community-based organization often support disaster response more quickly than an outside team during the first few hours?} \hfill [Tag: Role of community-based organizations in disaster response | Medium]
\begin{enumerate}[label=(\alph*),topsep=2pt,itemsep=1pt,partopsep=0pt,parsep=0pt]
    \item {It can ignore official warnings because local experience guarantees accurate predictions}
    \item {It can legally replace every specialized emergency service}
    \item {It always possesses more equipment than national response agencies}
    \item {It knows local hazards, vulnerable residents, routes, and resources}
\end{enumerate}
\vspace{2mm}

\noindent\textbf{Q40. During an earthquake response, an unverified social media post claims that a dam has collapsed and causes panic. What is the most responsible action for a news organization?} \hfill [Tag: Role of media in disaster response | Medium]
\begin{enumerate}[label=(\alph*),topsep=2pt,itemsep=1pt,partopsep=0pt,parsep=0pt]
    \item {Reveal the personal details of everyone who shared the post to discourage future rumors}
    \item {Verify the claim with authorized and independent reliable sources before reporting}
    \item {Broadcast the post immediately because rapid reporting is always the priority}
    \item {Repeat the claim while stating that its accuracy has not been checked}
\end{enumerate}
\vspace{2mm}

\noindent\textbf{Q41. Two districts experience earthquakes of equal magnitude. District X suffers far greater mortality because of unreinforced buildings, dense settlement, and nighttime occupancy. Which interpretation best explains the difference?} \hfill [Tag: Natural disasters | Hard]
\begin{enumerate}[label=(\alph*),topsep=2pt,itemsep=1pt,partopsep=0pt,parsep=0pt]
    \item {District X had greater exposure and vulnerability despite a similar hazard}
    \item {District X experienced a higher-magnitude hazard that instruments failed to detect}
    \item {District X experienced a natural hazard while the other district did not}
    \item {District X had lower resilience because earthquake frequency was necessarily higher}
\end{enumerate}
\vspace{2mm}

\noindent\textbf{Q42. An urban basin receives the same total rainfall in two storms. Storm P produces a much higher flood peak because rainfall is concentrated into one hour after extensive paving. Which mechanism most directly explains the higher peak?} \hfill [Tag: Water-related disasters: floods and drought | Hard]
\begin{enumerate}[label=(\alph*),topsep=2pt,itemsep=1pt,partopsep=0pt,parsep=0pt]
    \item {Increased infiltration and a longer groundwater residence time}
    \item {Reduced infiltration and a shorter time of concentration}
    \item {Reduced evaporation and a lower antecedent soil moisture level}
    \item {Increased channel storage and a flatter runoff hydrograph}
\end{enumerate}
\vspace{2mm}

\noindent\textbf{Q43. A region has near-normal annual rainfall, yet reservoirs, crops, and groundwater supplies fail because most rain falls in a few intense events followed by long dry spells. Which assessment is most defensible?} \hfill [Tag: Water-related disasters: floods and drought | Hard]
\begin{enumerate}[label=(\alph*),topsep=2pt,itemsep=1pt,partopsep=0pt,parsep=0pt]
    \item {Only meteorological drought can occur when annual rainfall is near normal}
    \item {Annual rainfall alone understates agricultural and hydrological drought}
    \item {Normal annual rainfall rules out every operational definition of drought}
    \item {Reservoir failure proves that the region experienced below-normal annual rainfall}
\end{enumerate}
\vspace{2mm}

\noindent\textbf{Q44. A severe cyclone is approaching a low-lying coast at high tide. Forecast uncertainty remains about the exact landfall point. Which action best applies risk-based evacuation planning?} \hfill [Tag: Air-related disasters: cyclones and storms | Hard]
\begin{enumerate}[label=(\alph*),topsep=2pt,itemsep=1pt,partopsep=0pt,parsep=0pt]
    \item {Evacuate surge-prone zones across the forecast uncertainty corridor}
    \item {Base evacuation solely on expected wind speed at the cyclone center}
    \item {Delay evacuation until the storm surge is directly observed onshore}
    \item {Evacuate only settlements located at the predicted landfall coordinate}
\end{enumerate}
\vspace{2mm}

\noindent\textbf{Q45. After the eye of a tropical cyclone passes over a coastal town, residents observe calm conditions and begin leaving shelters. What is the most appropriate official message?} \hfill [Tag: Air-related disasters: cyclones and storms | Hard]
\begin{enumerate}[label=(\alph*),topsep=2pt,itemsep=1pt,partopsep=0pt,parsep=0pt]
    \item {Remain sheltered because the eyewall may return with winds from the opposite direction}
    \item {Leave shelters because central pressure rises permanently once the eye arrives}
    \item {Open windows because equalizing pressure prevents most structural wind damage}
    \item {Move toward the coast because storm surge ends before the second eyewall}
\end{enumerate}
\vspace{2mm}

\noindent\textbf{Q46. A deep earthquake and a shallow earthquake have the same moment magnitude and occur at similar epicentral distances from a city. Which outcome is generally more likely, assuming comparable geology?} \hfill [Tag: Earth-related disasters: earthquakes, landslides, avalanches and volcanic eruptions | Hard]
\begin{enumerate}[label=(\alph*),topsep=2pt,itemsep=1pt,partopsep=0pt,parsep=0pt]
    \item {The shallow earthquake produces stronger local shaking because its rupture is closer}
    \item {Both earthquakes produce identical damage because epicentral distance is unchanged}
    \item {The deep earthquake produces stronger local shaking because magnitude rises with depth}
    \item {Both earthquakes produce identical shaking because moment magnitude fixes intensity}
\end{enumerate}
\vspace{2mm}

\noindent\textbf{Q47. Following a wildfire, intense rain is forecast over a steep valley containing settlements below burned slopes. Which intervention should receive immediate priority?} \hfill [Tag: Earth-related disasters: earthquakes, landslides, avalanches and volcanic eruptions | Hard]
\begin{enumerate}[label=(\alph*),topsep=2pt,itemsep=1pt,partopsep=0pt,parsep=0pt]
    \item {Plant deep-rooted trees before rainfall begins to stabilize every burned slope}
    \item {Issue debris-flow warnings using rainfall thresholds and evacuate drainage paths}
    \item {Excavate the slope toe to increase channel capacity beside the settlements}
    \item {Wait for visible slope cracks because debris flows require prior ground deformation}
\end{enumerate}
\vspace{2mm}

\noindent\textbf{Q48. Several avalanche victims are buried after a slab release. Rescuers have transceivers, probes, and shovels, while a second avalanche remains possible. Which sequence best balances rescue probability and responder safety?} \hfill [Tag: Earth-related disasters: earthquakes, landslides, avalanches and volcanic eruptions | Hard]
\begin{enumerate}[label=(\alph*),topsep=2pt,itemsep=1pt,partopsep=0pt,parsep=0pt]
    \item {Assess secondary risk, conduct a transceiver search, pinpoint by probing, then dig}
    \item {Wait for snowmelt, mark the crown line, then enter without electronic searching}
    \item {Begin random digging immediately, assess secondary risk, then activate transceivers}
    \item {Probe the entire debris field, move uphill, then search using transceivers}
\end{enumerate}
\vspace{2mm}

\noindent\textbf{Q49. A volcano shows increasing sulfur dioxide emissions, shallow seismic swarms, and ground inflation, but no eruption has begun. Which conclusion is scientifically appropriate?} \hfill [Tag: Earth-related disasters: earthquakes, landslides, avalanches and volcanic eruptions | Hard]
\begin{enumerate}[label=(\alph*),topsep=2pt,itemsep=1pt,partopsep=0pt,parsep=0pt]
    \item {Sulfur dioxide emissions exclude the possibility of a later phreatic eruption}
    \item {A major explosive eruption is certain within the next twenty-four hours}
    \item {Magma may be moving upward, but the timing and occurrence of eruption remain uncertain}
    \item {Ground inflation proves that pressure is declining within the magma system}
\end{enumerate}
\vspace{2mm}

\noindent\textbf{Q50. An earthquake ruptures industrial pipelines, causing a toxic release and fire near residential areas. How should this event be classified for planning purposes?} \hfill [Tag: Manmade disasters | Hard]
\begin{enumerate}[label=(\alph*),topsep=2pt,itemsep=1pt,partopsep=0pt,parsep=0pt]
    \item {A purely manmade disaster because industrial equipment caused every hazard}
    \item {A purely natural disaster because the initiating event was an earthquake}
    \item {A cascading disaster involving a natural trigger and technological consequences}
    \item {An unrelated pair of emergencies requiring no integrated command structure}
\end{enumerate}
\vspace{2mm}

\noindent\textbf{Q51. A radioactive plume is expected to pass over a community within thirty minutes, while evacuation would take two hours and create severe traffic congestion. What is the best immediate protective action?} \hfill [Tag: Nuclear disasters | Hard]
\begin{enumerate}[label=(\alph*),topsep=2pt,itemsep=1pt,partopsep=0pt,parsep=0pt]
    \item {Evacuate immediately through the projected plume regardless of travel time}
    \item {Distribute ordinary table salt as a substitute for official iodine prophylaxis}
    \item {Shelter indoors, close ventilation, and await monitored evacuation instructions}
    \item {Move residents outdoors so radioactive material does not accumulate around them}
\end{enumerate}
\vspace{2mm}

\noindent\textbf{Q52. Following a reactor release, officials distribute potassium iodide to an exposed population. Which statement correctly describes its protective role?} \hfill [Tag: Nuclear disasters | Hard]
\begin{enumerate}[label=(\alph*),topsep=2pt,itemsep=1pt,partopsep=0pt,parsep=0pt]
    \item {It blocks thyroid uptake of radioactive iodine but does not protect the whole body}
    \item {It removes all radionuclides already deposited in the lungs and digestive tract}
    \item {It replaces evacuation, sheltering, food controls, and environmental monitoring}
    \item {It neutralizes external gamma radiation from contaminated soil and buildings}
\end{enumerate}
\vspace{2mm}

\noindent\textbf{Q53. A chlorine tank ruptures near a town during calm nighttime conditions. The release forms a dense ground-level cloud. Which public instruction is most appropriate for people unable to evacuate safely?} \hfill [Tag: Chemical disasters | Hard]
\begin{enumerate}[label=(\alph*),topsep=2pt,itemsep=1pt,partopsep=0pt,parsep=0pt]
    \item {Move indoors to an upper level, seal openings, and stop outside-air ventilation}
    \item {Remain outdoors on low ground because chlorine rises rapidly when released}
    \item {Approach the leak from downwind while covering the face with a dry cloth}
    \item {Move into basements, open windows, and operate fans drawing outdoor air inward}
\end{enumerate}
\vspace{2mm}

\noindent\textbf{Q54. Hospitals report an unusual respiratory syndrome in geographically separated patients who attended the same indoor event. Human-to-human transmission is not yet confirmed. Which initial response is most appropriate?} \hfill [Tag: Biological disasters | Hard]
\begin{enumerate}[label=(\alph*),topsep=2pt,itemsep=1pt,partopsep=0pt,parsep=0pt]
    \item {Release patient identities publicly so untrained observers can locate additional cases}
    \item {Apply a case definition, isolate suspected cases, trace contacts, and investigate exposures}
    \item {Wait for proof of sustained transmission before beginning surveillance or infection control}
    \item {Treat every attendee as a confirmed case without laboratory or epidemiological assessment}
\end{enumerate}
\vspace{2mm}

\noindent\textbf{Q55. A tanker displaying a hazardous-material placard overturns in a tunnel, and an unidentified vapor is visible. What should the first arriving responder do before attempting rescue?} \hfill [Tag: Transport accidents | Hard]
\begin{enumerate}[label=(\alph*),topsep=2pt,itemsep=1pt,partopsep=0pt,parsep=0pt]
    \item {Wash the spill into tunnel drains to restore traffic and reduce surface contamination}
    \item {Approach from upwind, isolate the area, identify the substance remotely, and request specialists}
    \item {Ignite the vapor deliberately so its chemical identity becomes easier to determine}
    \item {Enter from downwind, inspect the tanker closely, and move casualties without protective equipment}
\end{enumerate}
\vspace{2mm}

\noindent\textbf{Q56. A flood affects several districts within one Indian state but does not exceed the state's response capacity. Which arrangement best reflects the national disaster management framework?} \hfill [Tag: National disaster management framework | Hard]
\begin{enumerate}[label=(\alph*),topsep=2pt,itemsep=1pt,partopsep=0pt,parsep=0pt]
    \item {The state leads through its authorities while district bodies execute local response plans}
    \item {The national authority automatically assumes direct command of every district operation}
    \item {Local bodies suspend their plans until Parliament declares a national disaster}
    \item {The armed forces permanently replace civil administration throughout the response}
\end{enumerate}
\vspace{2mm}

\noindent\textbf{Q57. A district disaster plan lists agencies and resources but assigns no trigger levels, reporting relationships, or authority for evacuation. Which failure is most likely during a fast-moving emergency?} \hfill [Tag: National disaster management framework | Hard]
\begin{enumerate}[label=(\alph*),topsep=2pt,itemsep=1pt,partopsep=0pt,parsep=0pt]
    \item {Automatic transfer of district powers to NGOs because evacuation authority is unstated}
    \item {Complete prevention of interagency communication because agency names are documented}
    \item {Excessive scientific certainty because resource inventories determine hazard probability}
    \item {Delayed decisions and duplicated actions because command and activation rules are unclear}
\end{enumerate}
\vspace{2mm}

\noindent\textbf{Q58. After a major earthquake, separate government departments issue conflicting casualty figures and independently dispatch teams to the same locations. Which corrective measure is most effective?} \hfill [Tag: Role of governmental agencies in disaster response | Hard]
\begin{enumerate}[label=(\alph*),topsep=2pt,itemsep=1pt,partopsep=0pt,parsep=0pt]
    \item {Suspend information exchange until all rescue operations have been completed}
    \item {Establish unified coordination, a common operating picture, and validated reporting channels}
    \item {Allow each department to publish estimates quickly without reconciling definitions}
    \item {Centralize every field decision with one remote official regardless of local conditions}
\end{enumerate}
\vspace{2mm}

\noindent\textbf{Q59. Several NGOs arrive after a cyclone and all propose distributing food in accessible neighborhoods, while remote settlements lack water and medical support. What is the best coordination approach?} \hfill [Tag: Role of NGOs in disaster response | Hard]
\begin{enumerate}[label=(\alph*),topsep=2pt,itemsep=1pt,partopsep=0pt,parsep=0pt]
    \item {Allocate every NGO the same neighborhood so beneficiaries can select among providers}
    \item {Permit independent distribution because competition always improves humanitarian coverage}
    \item {Use shared needs assessments and assign activities by location, sector, and capacity}
    \item {Exclude NGOs from response planning until government relief operations are entirely complete}
\end{enumerate}
\vspace{2mm}

\noindent\textbf{Q60. A community organization maintains a list of residents needing evacuation assistance. During a flash-flood warning, internet and mobile networks fail. Which preparedness measure would make the list most operationally useful?} \hfill [Tag: Role of community-based organizations in disaster response | Hard]
\begin{enumerate}[label=(\alph*),topsep=2pt,itemsep=1pt,partopsep=0pt,parsep=0pt]
    \item {Preassign nearby volunteers, backup communication methods, routes, and transport responsibilities}
    \item {Publish residents' medical details publicly so any passerby can make evacuation decisions}
    \item {Wait for external responders because local volunteers should never support evacuation}
    \item {Store the list only in a remote cloud system to maximize confidentiality during outages}
\end{enumerate}
\vspace{2mm}


\section{Section B: Comprehensive Subjective Problem with Detailed Analytical Rubric}

\begin{subjectiveq}{5}{1}
\textbf{Integrated Disaster Risk Assessment and Flash Flood Mitigation:}\\The United Nations Office for Disaster Risk Reduction (UNDRR) and the Sendai Framework define disaster risk through the multi-variable function:\begin{equation}\text{Disaster Risk} = \frac{\text{Hazard} \times \text{Vulnerability} \times \text{Exposure}}{\text{Capacity and Coping Mechanisms}}\end{equation}(a) Define each of the four components in the context of mountain river basins subject to Glacial Lake Outburst Floods (GLOFs) and cloudburst events (e.g. 2013 Kedarnath, 2021 Chamoli events in Uttarakhand, India).\\(b) A Himalayan pilgrimage town of $25{,}000$ permanent residents receives an influx of $60{,}000$ floating pilgrim tourists during the monsoon month of July. Unregulated concrete construction encroaches within $15\text{ meters}$ of the river high-flood line. Early warning Doppler weather radar coverage is absent, and the single mountain access road is prone to debris landslides. Conduct a risk audit by scoring each parameter on a normalized index ($1$ to $10$) and calculate the composite risk quotient.\\(c) Design a multi-tiered disaster management protocol spanning: (i) Pre-disaster prevention and structural zoning, (ii) Early warning telemetry, (iii) Real-time search and rescue (NDRF, SDRF, Army), and (iv) Post-disaster Sendai ``Build Back Better'' rehabilitation.
\end{subjectiveq}

\begin{solbox}
\textbf{Part (a): Conceptual Definition of Risk Elements:}\\\begin{itemize}[leftmargin=1.5em]  \item \textbf{Hazard ($H$):} The intrinsic physical frequency, intensity, and probability of an extreme geomorphic event (e.g., $150\text{ mm/hr}$ cloudburst or moraine dam collapse yielding a $2{,}500\text{ m}^3/\text{s}$ flash flood wave).  \item \textbf{Exposure ($E$):} The spatial presence of human lives, critical infrastructure, bridges, roads, hotels, and pilgrims situated within the active hydrologic flood inundation zone.  \item \textbf{Vulnerability ($V$):} Structural susceptibility and systemic fragility, including unreinforced masonry on uncompacted glacial debris, lack of emergency evacuation pathways, and high proportion of elderly tourists.  \item \textbf{Capacity ($C$):} The collective institutional, logistical, and technical capability to forecast, communicate early warnings, shelter, evacuate, and deliver medical emergency response.\end{itemize}\textbf{Part (b): Quantitative Disaster Risk Index Audit:}\\Assigning normalized baseline scores ($1 = \text{lowest}, 10 = \text{critical}$):\begin{align*}\text{Hazard Score ($H$)} &= 9.0\text{ (Severe alpine terrain, high cloudburst probability, steep hydraulic gradient)} \\\text{Exposure Score ($E$)} &= 8.5\text{ ($85{,}000$ individuals situated directly on the riparian floodway)} \\\text{Vulnerability Score ($V$)} &= 9.0\text{ (Weak foundations, single fragile egress route, zero flood barriers)} \\\text{Capacity Score ($C$)} &= 2.5\text{ (No Doppler radar, limited satellite telemetry, remote response units)}.\end{align*}Evaluating the composite Risk Quotient:\begin{equation*}\text{Risk} = \frac{H \times E \times V}{C} = \frac{9.0 \times 8.5 \times 9.0}{2.5} = \frac{688.5}{2.5} = 275.4.\end{equation*}If community disaster management capacity is upgraded by deploying automated weather stations (AWS), SMS early warnings, dedicated evacuation shelters, and satellite comms ($C \to 8.0$), while enforcing strict river corridor setback zones ($E \to 3.0$):\begin{equation*}\text{Mitigated Risk} = \frac{9.0 \times 3.0 \times 4.0}{8.0} = \frac{108.0}{8.0} = 13.5\text{ (a $95.1\%$ reduction in disaster risk)}.\end{equation*}\textbf{Part (c): Multi-Tiered Disaster Framework:}\\\begin{itemize}[leftmargin=1.5em]  \item \textbf{Phase 1 (Prevention \& Preparedness):} Enactment of River Regulation Zones (RRZ); geotechnical landslide stabilization using soil nailing, gabion retaining walls, and reforestation; community mock drills.  \item \textbf{Phase 2 (Telemetry \& Warning):} Automated lake water level sensors linked via INSAT uplinks providing minimum $45\text{-minute}$ advance evacuation siren warnings.  \item \textbf{Phase 3 (Incident Command \& Rescue):} Rapid mobilization of National Disaster Response Force (NDRF) battalions, Indian Air Force helipads, and satellite phones.  \item \textbf{Phase 4 (Build Back Better):} Prohibiting reconstruction on ancient landslide fans; building seismically resilient green hospitals and public infrastructure.\end{itemize}
\end{solbox}

\begin{examinertip}
\textbf{Marking Scheme Breakdown (Total: 10 Marks):}\\
\textbullet\ \textbf{3 Marks:} Clear formulation of governing environmental identities, equations, and thermodynamic/constitutional foundations.\\
\textbullet\ \textbf{4 Marks:} Rigorous algebraic/numerical substitution, dimensional consistency, and unit handling.\\
\textbullet\ \textbf{3 Marks:} Critical evaluation of environmental implications, regulatory criteria (CPCB/UNDRR/IUCN), and policy mitigations.
\end{examinertip}
